\documentclass[10pt,a4paper]{amsart}
\usepackage[margin=19mm]{geometry}
\usepackage{amsmath,amssymb,mathtools}
\usepackage[scr=rsfs,cal=euler]{mathalfa}
\usepackage{xfrac}
\usepackage[unicode,hidelinks,bookmarks=false]{hyperref}
\newcommand{\ie}{i.e.\ }
\newcommand{\cf}{cf.\ }
\newcommand{\resp}{respectively\ }
\newcommand{\Subsection}{\subsection}
\providecommand{\nobreakdash}{\nobreak\hbox{-}}
\setlength{\emergencystretch}{2em}
\multlinegap0pt
\usepackage{bm}
\usepackage{bbm}
\usepackage{dsfont}
\usepackage{xspace}
\usepackage{cancel}
\usepackage{mathtools}
\usepackage{amsthm}
\makeatletter
\@ifundefined{dummy}{\newtheorem{dummy}{Dummy}}{}
\@ifundefined{lem}{\newtheorem{lem}[dummy]{Lemma}}{}
\makeatother
\title[The point insertion technique and open $r$-spin theories I: ]{The point insertion technique and open $r$-spin theories I: moduli and orientation}
\author{Ran J. Tessler, Yizhen Zhao}
\date{Source: arXiv:2310.13185v4 (CC BY 4.0)}
\begin{document}
\maketitle

\noindent\emph{Source and licence.} The point insertion technique and open $r$-spin theories I: moduli and orientation. Version arXiv:2310.13185v4 (CC BY 4.0), distributed under the Creative Commons Attribution 4.0 International licence, as recorded in the arXiv licence metadata for this exact version and shown on the abstract page https://arxiv.org/abs/2310.13185v4. Full rights notice: licenses/arxiv-2310.13185-CC-BY-4.0.txt.\par\smallskip

\noindent\textit{Editorial scope.} Counted unit: the Lem whose source label is \texttt{None} in the article (source0.tex lines 1456--1458, source label \texttt{None}), delivered together with the article's own complete original proof of it (source0.tex lines 1459--1528, one \texttt{proof} environment of 70 lines). No mathematics is written, repaired or re-derived: statement and proof are the article's own text line for line. Required context retained uncounted: eq change of orientation of witten when cross a legal in g 1 1st bdry (lines 1447--1449); eq change of orientation of witten when cross a legal in g 1 2nd bdry (lines 1451--1453). Every definition, display and label of those lines is retained unchanged. Omitted from this package and not redistributed: 1--1446, 1450--1450, 1454--1455, 1529--3011. Nothing inside a retained excerpt is omitted or altered: every retained body is delivered exactly as the article's lines are written, so no omission receipt of any kind accompanies this package. Citations: the retained excerpts carry no citation command at all, so no bibliography entry of the article is transcribed and no reference is redistributed. Reference adaptation: none was needed and none was made; every reference inside the delivered unit resolves inside it, so no unresolved reference or citation is produced. Printed slips kept literal: no printed word, symbol, hypothesis or number of the counted statement or of its proof was altered. Layout and class: the article's own class is \texttt{article} with packages including inputenc, geometry, amsmath, amsthm, amssymb, amscd, dsfont, bm, epic, accents, graphicx, xy, tikz, multicol; the wrapper is the lane's pinned \texttt{amsart} editorial wrapper at 10pt on A4, which restates the theorem-like environments the retained text needs and adds the article's own macro definitions that the retained text needs (none), with extra wrapper packages (bm, bbm, dsfont, xspace, cancel, mathtools). The delivered numbering is the unit's own. Assets: no retained span contains a figure environment, a graphics-inclusion command, a PDF-page inclusion, a table or a TikZ picture; no image file is copied into this package, the package declares no asset dependency and no image file is redistributed with it. Licence: version arXiv:2310.13185v4 of this article is distributed under the Creative Commons Attribution 4.0 International licence, as recorded in the arXiv licence metadata for this exact version and shown on the abstract page https://arxiv.org/abs/2310.13185v4. A full rights notice is delivered at licenses/arxiv-2310.13185-CC-BY-4.0.txt. Characters: every retained excerpt is pure ASCII, so the delivered engine, XeLaTeX with its default Latin Modern text font, reports no missing character.
\begin{equation}\label{eq change of orientation of witten when cross a legal in g 1 1st bdry}
\mathfrak o(p_\alpha,p_\beta,N_\alpha,N_\beta)=(-1)^{N_\alpha}\mathfrak o(p'_\alpha,p_\beta,N_\alpha,N_\beta);
\end{equation} 

\begin{equation}\label{eq change of orientation of witten when cross a legal in g 1 2nd bdry}
\mathfrak o(p_\alpha,p_\beta,N_\alpha,N_\beta)=(-1)^{N_\beta}\mathfrak o(p_\alpha,p'_\beta,N_\alpha,N_\beta).
\end{equation}

\begin{lem}
    If $N_\alpha \ge 2$, we have  $\mathfrak o(p_\alpha,p_\beta,N_\alpha,N_\beta)=\mathfrak o(p_\alpha,p_\beta,N_\alpha-2,N_\beta+2)$.
\end{lem}

\begin{proof}
    We only need to check it for the fibre of $\mathcal W$ over a fixed cylinder $(C,\Sigma)$. According to \eqref{eq change of orientation of witten when cross a legal in g 1 1st bdry} and \eqref{eq change of orientation of witten when cross a legal in g 1 2nd bdry}, we only need to prove it for a specific choice of $p_\alpha,p_\beta$.

    For each cylinder $(C,\Sigma)$ in $\mathcal M^*$, we choice a coordinate system and represent $C$ as $\mathbb C=\mathbb R_x\times R_y$ quotiented by the lattice generated by $(1,0)$ and $(0,2\tau)$, and the fundamental domain of $\Sigma$ is $\{0\le x < 1\}\times \{0\le y < \tau\}$. We assume $(0,0)$ and $(0,\tau)$ are now boundary markings, otherwise we can shit the $x$-coordinator a little bit.

    We fix $p_1=(0,0)$ and $p_2=(0,\tau).$ If we choice $q_{\alpha,1},q_{\alpha,2},\dots,q_{\alpha,{N_\alpha}}$ and $q_{\beta,1},q_{\beta,2},\dots,q_{\beta,{N_\beta}}$ to define $\mathfrak o(p_\alpha,p_\beta,N_\alpha,N_\beta)$ as 
    $$
    s_{\alpha,1}\dots \wedge s_{\alpha,N_\alpha}\wedge s_{\beta,1}\wedge \dots \wedge s_{\beta,N_\beta},
    $$
    then we have $q_{\alpha,i}=(x_{\alpha,i},0)$ for each $0\le i \le N_\alpha$, where $0< x_{\alpha,1}<x_{\alpha,2}<\dots<x_{\alpha,N_\alpha}<1$; we also have $q_{\beta,i}=(x_{\beta,i},\tau)$ for each $1\le i \le N_\beta$, where $1> x_{\beta,1}>x_{\beta,2}>\dots>x_{\beta,N_\beta}>0$. We write $\check{q}_{\alpha,i}=(\check{x}_{\alpha,i},\tau)$ and $\check{q}_{\beta,i}=(\check{x}_{\beta,i},0)$.
    Since $s_{j,i}$ are sections of a fixed bundle $\lvert J \rvert$, the sum of (the $x$-coordinators) its zeros 
    \begin{equation}\label{eq sum of zero of x coordinator of a section}
        X:=\sum_{k=1}^{N_\alpha} x_{\alpha,k}+  \sum_{k=1}^{N_\beta} x_{\beta,k} -x_{j,i} + \check{x}_{j,i}
    \end{equation}
    is a constant modulo $\mathbb Z$ for a fixed coordinate system. We can assume $0< X < 1$, otherwise we can shift the $x$-coordinator in our coordinate system.

    We take  $q_{\alpha,1},q_{\alpha,2},\dots,q_{\alpha,{N_\alpha}}$ and $q_{\beta,1},q_{\beta,2},\dots,q_{\beta,{N_\beta}}$ in a way that 
    $$0< x_{\alpha,1}<x_{\alpha,2}<\dots<x_{\alpha,N_\alpha}<\frac{\operatorname{min}\{X,1-X\}}{2N}$$ 
    and 
    $$1> x_{\beta,1}>x_{\beta,2}>\dots>x_{\beta,N_\beta}>1-\frac{\operatorname{min}\{X,1-X\}}{2N},$$ 
    are require then do not touch the boundary markings. Then by \eqref{eq sum of zero of x coordinator of a section} we have 
    $$1> x_{\beta,1}>x_{\beta,2}>\dots>x_{\beta,N_\beta}> \check{x}_{\alpha,N_\alpha}> \check{x}_{\alpha,N_\alpha-1}>0.$$

    Now we choice $q'_{\alpha,1},q'_{\alpha,2},\dots,q'_{\alpha,{N_\alpha-1}}$ and $q'_{\beta,1},q'_{\beta,2},\dots,q'_{\beta,{N_\beta}},q'_{\beta,{N_\beta+1}},q'_{\beta,{N_\beta+2}}$ to define the orientation $\mathfrak o(p_\alpha,p_\beta,N_\alpha-2,N_\beta+2)$ as 
    $$
    s'_{\alpha,1}\dots \wedge s'_{\alpha,N_\alpha-2}\wedge s'_{\beta,1}\wedge \dots \wedge s'_{\beta,N_\beta+2}.
    $$
    We take $q'_{\alpha,i}=q_{\alpha,i}$ for $1\le i \le N_\alpha-2$, $q'_{\beta,i}=q_{\beta,i}$ for $1\le i \le N_\beta$,  $q'_{\beta,N_\beta+1}=\check{q}_{\alpha,N_\alpha}$ and $q'_{\beta,N_\beta+2}=\check{q}_{\alpha,N_\alpha-1}$. 
    Under this choice, the sections $s_{j,i}$ and $s'_{j,i}$ satisfy the following properties.
    \begin{itemize}
        \item For any $1\le i \le N_\alpha-2$, both $s_{\alpha,i}$ and $s'_{\alpha,i}$ vanish when evaluated at $\{q_{\alpha,1},\dots,q_{\alpha,N_\alpha},q_{\beta,1},\dots,q_{\beta,N_\beta}\}\setminus\{q_{\alpha,i}\}$. Moreover, $s_{\alpha,i}(q_{\alpha,i})$ and $s'_{\alpha,i}(q_{\alpha,i})$ have the same sign with respect to the grading: $s_{\alpha,i}$ and $s'_{\alpha,i}$ have the same number of zeros on the arc from $p_\alpha$ to $q_{\alpha,i}$, and both $s_{\alpha,i}(p_\alpha)$ and $s'_{\alpha,i}(p_\alpha)$ are negative by definition. Therefore, for each $0\le t \le 1$, the section
        $$
        s^t_{\alpha,i}:=(1-t)s_{\alpha,i}+ts'_{\alpha,i}
        $$
        vanishes when evaluated at $\{q_{\alpha,1},\dots,q_{\alpha,N_\alpha},q_{\beta,1},\dots,q_{\beta,N_\beta}\}\setminus\{q_{\alpha,i}\}$, and does not vanish when  evaluated at $q_{\alpha,i}$.
        
        \item For any $1\le i \le N_\beta$, similar to the above item, for each for each $0\le t \le 1$, the section
        $$
        s^t_{\beta,i}:=(1-t)s_{\beta,i}+ts'_{\beta,i}
        $$
        vanishes when evaluated at $\{q_{\alpha,1},\dots,q_{\alpha,N_\alpha},q_{\beta,1},\dots,q_{\beta,N_\beta}\}\setminus\{q_{\beta,i}\}$, and does not vanish when  evaluated at $q_{\beta,i}$.
        
        \item The zeros of $s_{\alpha,N_\alpha-1}$ and $s'_{\beta,N_\beta+1}$ are the same: at $q_{\alpha,i}$ for $1\le i \le N_\alpha-2$, $q_{\beta,i}$ for $1\le i \le N_\beta$, $q_{\alpha,N_\alpha}$ and $\check{q}_{1,N_\alpha-1}$. Thus $\frac{s_{\alpha,N_\alpha-1}}{s'_{\beta,N_\beta+1}}$ is a non-zero real number  $\frac{s_{\alpha,N_\alpha-1}(p_\beta)}{s'_{\beta,N_\beta+1}(p_\beta)}$. Similarly, $\frac{s_{\alpha,N_\alpha}}{s'_{\beta,N_\beta+2}}$ is a non-zero real number  $\frac{s_{\alpha,N_\alpha}(p_\beta)}{s'_{\beta,N_\beta+2}(p_\beta)}$. Note that $s'_{\beta,N_\beta+1}(p_\beta)$ and $s'_{\beta,N_\beta+2}(p_\beta)$ are always negative by definition. 
        Moreover, we can show that $s_{\alpha,N_\alpha}(p_\beta)$ and $s_{\alpha,N_\alpha-1}(p_\beta)$ always have the same sign. In fact, we can construct a homotopy $s^\nu$ between $s_{\alpha,N_\alpha}$ and $s_{\alpha,N_\alpha-1}$ with a homotopy parameter $0\le \nu \le 1$: the zeros of $s^\nu$ are at $q_{\alpha,i}$ for $1\le i \le N_\alpha-2$, $q_{\beta,i}$ for $1\le i \le N_\beta$, $\nu q_{\alpha,N_\alpha}+(1-\nu)q_{\alpha,N_\alpha-1}$ and $\nu\check{q}_{\alpha,N_\alpha-1}+(1-\nu)\check{q}_{\alpha,N_\alpha}$. 
        Then for any $0\le \nu \le 1$, $s^\nu$ does not vanish at $p_\beta$, which means $s_{\alpha,N_\alpha}(p_\beta)=s^{\nu=0}(p_\beta)$ and $s_{\alpha,N_\alpha-1}(p_\beta)=s^{\nu=1}(p_\beta)$ have the same sign. Therefore up to a positive scaling we have
        $$
        s_{\alpha,N_\alpha-1}\wedge s_{\alpha,N_\alpha}=s'_{\beta,N_\beta+1} \wedge s'_{\beta,N_\beta+2}.
        $$
    \end{itemize}
    Then the family of wedge products (over $0\le t \le 1$)
    $$
    s^t_{\alpha,1}\dots \wedge s^t_{\alpha,N_\alpha-2}\wedge s_{\alpha,N_\alpha-1}\wedge s_{\alpha,N_\alpha}\wedge s^t_{\beta,1}\wedge \dots \wedge s^t_{\beta,N_\beta},
    $$
    homotopes between 
    $$\mathfrak o(p_\alpha,p_\beta,N_\alpha,N_\beta)= s_{\alpha,1}\dots \wedge s_{\alpha,N_\alpha-2}\wedge s_{\alpha,N_\alpha-1}\wedge s_{\alpha,N_\alpha}\wedge s_{\beta,1}\wedge \dots \wedge s_{\beta,N_\beta}$$
    and 
    \begin{equation*}
    \begin{split}
        \mathfrak o(p_\alpha,p_\beta,N_\alpha-2,N_\beta+2)=& s'_{\alpha,1}\dots \wedge s'_{\alpha,N_\alpha-2}\wedge s'_{\beta,1}\wedge \dots \wedge s'_{\beta,N_\beta} \wedge s'_{\beta,N_\beta+1}\wedge s'_{\beta,N_\beta+2}\\
        =&s'_{\alpha,1}\dots \wedge s'_{\alpha,N_\alpha-2}\wedge s'_{\beta,1}\wedge \dots \wedge s'_{\beta,N_\beta} \wedge s_{\alpha,N_\alpha-1}\wedge s_{\alpha,N_\alpha}\\
        =&s'_{\alpha,1}\dots \wedge s'_{\alpha,N_\alpha-2} \wedge s_{\alpha,N_\alpha-1}\wedge s_{\alpha,N_\alpha}\wedge s'_{\beta,1}\wedge \dots \wedge s'_{\beta,N_\beta}.
    \end{split}
     \end{equation*}
    It remains to show that for any $0\le t \le 1$, the sections $s^t_{\alpha,1},\dots, s^t_{\alpha,N_\alpha-2}$, $s^t_{\beta,1}, \dots , s^t_{\beta,N_\beta}$, $ s_{\alpha,N_\alpha-1},$ $s_{\alpha,N_\alpha}$ are linearly independent. Assuming a linear combination of them are zero, \textit{i.e.,}
    $$
    \lambda^t_{\alpha,1}s^t_{\alpha,1}\dots + \lambda^t_{\alpha,N_\alpha-2}s^t_{\alpha,N_\alpha-2}+ \lambda_{\alpha,N_\alpha-1} s_{\alpha,N_\alpha-1}+ \lambda_{\alpha,N_\alpha} s_{\alpha,N_\alpha}+ \lambda^t_{\beta,1} s^t_{\beta,1}+ \dots + \lambda^t_{\beta,N_\beta} s^t_{\beta,N_\beta}=0,
    $$
    the evaluation at $q_{\alpha,i}$ forces $\lambda^t_{\alpha,i}=0$ for all $1\le i \le N_\alpha-2$ since $s^t_{\alpha,i}$ is the only one dose not vanish at $q_{\alpha,i}$. Similarly we have $\lambda^t_{\beta,i}=0$ for all $1\le i \le N_\beta$. Then the only terms left are 
    $$\lambda_{\alpha,N_\alpha-1} s_{\alpha,N_\alpha-1}+ \lambda_{\alpha,N_\alpha} s_{\alpha,N_\alpha}=0,$$ this forces $\lambda_{\alpha,N_\alpha-1}=\lambda_{\alpha,N_\alpha}=0$ since $s_{\alpha,N_\alpha-1}$ and $s_{\alpha,N_\alpha}$ have different zeros.
\end{proof}

\end{document}
